\documentclass[a4paper,11pt]{article}
\usepackage[margin=2.5cm]{geometry}
\usepackage{booktabs}
\usepackage{parskip}
\usepackage{makecell}

\title{\textbf{Lab 6 -- UVM Testbench}\\[0.4em]
       \large Existing Testbench Structure}
\author{Dan Hagen}
\date{May 2026}

\begin{document}

\maketitle

\begin{table}[h]
\centering
\begin{tabular}{llll}
  \toprule
  Block & Lines & Role & UVM equivalent \\
  \midrule
  Parameters \& signals & 1--27    & Declare DUT ports directly       & \texttt{mem\_ctrl\_if} \\ \hline
  DUT instantiation     & 30--34   & Connect TB to DUT                & (unchanged) \\ \hline
  Transaction type      & 36--50   & Combined M0+M1 stimulus packet   & \texttt{ctrl\_transaction} \\ \hline
  Mailboxes \& flags    & 52--65   & Inter-process communication      & \makecell[l]{\texttt{ctrl\_sequencer}\\+ objections} \\ \hline
  Clock \& Reset        & 67--77   & Timing infrastructure            & Top-level module \\ \hline
  Generator (directed)  & 89--119  & Deterministic M0/M1 stimulus     & \texttt{ctrl\_directed\_seq} \\ \hline
  Generator (contention)& 122--137 & Simultaneous M0+M1 stimulus      & \texttt{ctrl\_contention\_seq} \\ \hline
  Generator (random)    & 140--156 & Randomised mixed stimulus        & \texttt{ctrl\_random\_seq} \\ \hline
  Driver                & 172--233 & Drive DUT pins from mailbox      & \texttt{ctrl\_driver} \\ \hline
  Checker               & 236--284 & Golden model + result comparison & \texttt{ctrl\_scoreboard} \\ \hline
  Monitor               & 289--349 & Sample DUT signals, build events & \texttt{ctrl\_monitor} \\ \hline
  Main control          & 354--397 & \makecell[l]{Launch tasks,\\wait, \texttt{\$finish}} & \makecell[l]{\texttt{ctrl\_directed\_test}\\\texttt{ctrl\_contention\_test}\\\texttt{ctrl\_random\_test}\\\texttt{ctrl\_full\_test}} \\
  \bottomrule
\end{tabular}
\caption{Structure of \texttt{tb\_simple\_mem\_ctrl.sv} and its UVM mapping.}
\end{table}

\newpage
\section*{UVM Testbench Diagram}

\begin{verbatim}
uvm_test_top (ctrl_directed_test | ctrl_contention_test |
              ctrl_random_test   | ctrl_full_test)
|
+-- ctrl_env
    |
    +-- ctrl_agent
    |   |
    |   +-- ctrl_sequencer <-------- ctrl_directed_seq
    |   |         |                  ctrl_contention_seq
    |   |    TLM handshake           ctrl_random_seq
    |   |         |
    |   +-- ctrl_driver ----drive---> mem_ctrl_if ---> DUT (simple_mem_ctrl)
    |   |                                                       |
    |   +-- ctrl_monitor <--sample-- mem_ctrl_if <-------------+
    |         |
    |    analysis_port
    |         |
    +-- ctrl_scoreboard
\end{verbatim}


\section*{Ai usage}




Work alognside it  as a guid and help to find relvant parts in slides + some background info

Example : 


Ok lets move on to the monitor. For now just give relvant line numbers in the original tb, the     
relevant slide pages . And super short idea of what its purpose/fucntionality is                     
                   

Response:
⏺ Original TB: lines 289-349 (tb_monitor)                                                            
                                                                                                     
  Slides: W08 slide 18 (uvm_monitor role), W08 slide 31 (FIFO monitor implementation), W09 slide 4   
  (communication/connect phases), W09 slide 30 (analysis ports)                                      
                                                                                                     
  Purpose: Passively samples DUT signals every clock cycle, packages what it observes into a         
  ctrl_transaction, and broadcasts it via uvm_analysis_port to anyone listening (scoreboard). It   
  never drives anything — read-only. 




Response:

  Ok again for top module give me the relvant slides and info to try        
  implmenting it my slefe. Also  what does super.run_phase(phase) do ? Why do 
 the slides implement it in their examples ?                                
                                                                            
⏺ super.run_phase(phase): It calls the parent class's run_phase. In         
  uvm_test, that's an empty method — does nothing. The slides include it as
  a habit/convention to be safe in case you ever extend a test from another 
  test (where the parent might have real logic in run_phase). For your tests
   extending directly from uvm_test, it's optional. Up to you whether to    
  keep it for consistency with the slides or remove it.
                          
  ---                                            
  Top module:
             
  - Slides: W09 slides 5-6 (top module example with interface, DUT,
  config_db, run_test).                                                     
  - Original TB: Lines 67-77 (clock + reset), lines 31-34 (DUT
  instantiation), lines 356-364 (signal init).                              
                             


Repatative tasks: 

(after doing 2/3 display to uvm style):
Switch the rest of display to uvm style 


Linting + spelling + naming mismatches:

Clean up code (set in .md file that this is doesnt touch logic)

  
\end{document}
